% To be reformatted in the journal style upon acceptance.
\documentclass[11pt]{amsart}

\usepackage{amsmath,amssymb}
\usepackage[hidelinks]{hyperref}

\newtheorem{thm}{Theorem}[section]
\newtheorem{lem}[thm]{Lemma}
\newtheorem{cor}[thm]{Corollary}
\newtheorem{prop}[thm]{Proposition}
\theoremstyle{definition}
\newtheorem{defi}[thm]{Definition}
\newtheorem{rem}[thm]{Remark}
\newtheorem{obs}[thm]{Observation}

\newcommand{\I}{\mathbb{I}}
\newcommand{\DM}[1]{\mathrm{DM}(#1)}
\newcommand{\aconv}{\equiv_{\mathrm{alg}}}
\newcommand{\nconv}{\not\equiv}
\newcommand{\propeq}{\equiv_{\mathrm{prop}}}
\newcommand{\refl}{\mathsf{refl}}
\newcommand{\Path}{\mathsf{Path}}
\newcommand{\U}{\mathcal{U}}
\newcommand{\hcomp}{\mathsf{hcomp}}
\newcommand{\ua}{\mathsf{ua}}
\newcommand{\transp}{\mathsf{transp}}
\newcommand{\cng}{\mathsf{cong}}
\newcommand{\trns}{\mathbin{\cdot}}
\newcommand{\Ftwo}{F_2}
\newcommand{\SL}{\mathrm{SL}_2(\mathbb{Z})}
\newcommand{\RP}{\mathbb{RP}^2}
\newcommand{\Sone}{S^1}

\begin{document}

\title[Word problems across the strict--weak boundary]{Word problems
  across the strict--weak boundary\\ in cubical type theory}

\author{Ryota Shibusawa}
\address{Daiichi Institute of Technology, Japan}
\email{r-shibusawa@daiichi-koudai.ac.jp}

\subjclass[2020]{03D40, 03B38, 18N45, 20F10}
\keywords{cubical type theory, word problem, decidability, cubical
  computads, covering spaces, univalence, free groups, group
  presentations}

\begin{abstract}
  Over a cellular context --- the context of a cubical computad,
  freely listing cell variables over a De Morgan interval theory ---
  two word problems coexist: whether two definable cells are equal
  \emph{strictly} (by the decidable algorithmic equality of a cubical
  kernel) and whether they are equal \emph{up to homotopy} (by a
  propositional identification).  We show that the boundary between
  the strict and the weak layer is exactly the boundary between
  decidability and undecidability.  The strict word problem is
  decidable over every cellular context, in every dimension.  The
  propositional word problem is decidable over free graph contexts
  (dimension~$1$), where we make both halves of the decision
  procedure computational: identifications are compiled from generic
  groupoid laws, and separations are computed by a
  univalent transport invariant with a \emph{proof-free carrier} ---
  the faithful representation $\Ftwo \hookrightarrow \SL$ of Sanov,
  which evaluates at machine speed (under a millisecond per letter)
  where the reduced-word invariant is impractical.  From
  dimension~$2$ on, the propositional problem is undecidable:
  finitely presented groups embed as relator cells, and both
  directions of the reduction are exhibited inside the kernel ---
  relator applications compile to explicit path constructions over
  the generic presentation context, and the presentation complex of
  $\mathbb{Z}/2$ is realized internally as $\RP$, where a
  univalent double cover certifies that the presentation is realized
  faithfully ($x^2 = 1$ and $x \neq 1$, both machine-checked).  All
  positive claims are anchored in a kernel-checked artifact; the
  measured computational frontier of the kernel itself is reported.
\end{abstract}

\maketitle

\section{Introduction}
\label{sec:intro}

A cubical proof assistant maintains two equalities.  The
\emph{strict} (definitional, algorithmic) equality $\aconv$ is the
one its kernel decides silently at every type check; the \emph{weak}
(propositional) equality $\propeq$ is the one its users prove, path
by path.  A companion line of work
\cite{Const,Sphere,Realization,Convexity} has measured the gap
between the two layers over \emph{cellular contexts} --- contexts
$\Gamma$ that freely list cell variables ($A : \U$, points, paths,
squares, \dots), the type-theoretic incarnation of cubical computads
--- and found the strict layer to be small, combinatorial, and
well-organized: strict cells are classified by reparametrizations
along the free De Morgan algebra, the strict layer of a generic
$n$-cell is a wedge of De Morgan spheres, cellular contexts realize
their cube presheaves on the nose, and the strict fibers over a
boundary are strictly contractible.

This paper asks the decision-theoretic question: \emph{given two
definable cells over a cellular context, can a machine decide whether
they are equal?}  The answer splits along the same seam:

\begin{center}
\begin{tabular}{l|l|l}
  & dimension $1$ (free) & dimension $\geq 2$ (with relators)\\
  \hline
  strict $\aconv$ & decidable & \textbf{decidable}\\
  weak $\propeq$ & decidable & \textbf{undecidable}
\end{tabular}
\end{center}

\noindent The strict word problem is decidable over \emph{every}
cellular context (Theorem~\ref{thm:strict}): the kernel's
algorithmic equality is a decision procedure, and the realization
theorem of \cite{Realization} says exactly that it decides the
intended combinatorial equality.  The propositional word problem is
decidable in dimension~$1$ (Theorem~\ref{thm:free}) and undecidable
from dimension~$2$ on (Theorem~\ref{thm:undec}): group presentations
embed as relator cells and the word problem of a finitely presented
group reduces to $\propeq$, so Novikov--Boone applies
\cite{Novikov,Boone}.  \emph{The strict--weak boundary is the
decidable--undecidable boundary.}  This gives a quantitative reading
of what weakening buys and what it costs: relator cells create the
identifications that make homotopy theory expressive, and those same
identifications destroy decidability --- while the strict layer,
which by the no-go theorems of \cite{Const,Convexity} can never
absorb a relator, remains decidable throughout.

The mathematical content of the boundary statement is classical:
normal forms for free groupoids on one side, Novikov--Boone on the
other.  The contribution of this paper is to make the statement
\emph{computational} in a live cubical kernel, in four steps.

\begin{enumerate}
\item \emph{Identifications compile.}  Over the free context, the
  identification half of the dimension-$1$ decision procedure is a
  compiler from reduction sequences of words to explicit path
  composites of generic groupoid laws, each machine-checked once and
  for all (\S\ref{sec:free}).  Over a presentation context, relator
  applications compile the same way: the group-theoretic consequence
  $x^4 = 1$ of the single relation $x^2 = 1$ becomes an explicit
  path constructed from two relator uses and one unit law, checked
  over the \emph{generic} presentation context
  (\S\ref{sec:relators}).
\item \emph{Separations compute.}  The separation half needs
  invariants that evaluate.  Abelian invariants (winding numbers)
  are cheap but incomplete.  We introduce a transport invariant with
  a \emph{proof-free carrier}: the faithful Sanov representation
  $\Ftwo \hookrightarrow \SL$, implemented as a univalent cover of
  the figure eight whose fibers are $\mathbb{Z}^4$ and whose
  monodromies are elementary integer-arithmetic equivalences.
  Replacing a proof-carrying carrier (reduced words as a subset
  type) by a flat one moves the per-letter cost from ``infeasible''
  to under a millisecond, and faithfulness makes the invariant
  complete for $\Ftwo$ (\S\ref{sec:sanov}).
\item \emph{Presentations realize.}  For the undecidability
  transfer, the presentation complex is built \emph{inside} the
  theory: for $\langle x \mid x^2 \rangle$ it is the mapping cone of
  the degree-$2$ map, i.e.\ $\RP$, and the relator cell folds into a
  machine-checked identification $x \trns x \propeq \refl$ --- one
  $\hcomp$ square capped at the cone point.  Generic relator
  derivations then transfer by instantiation alone
  (\S\ref{sec:present}).
\item \emph{Faithfulness certifies.}  A univalent double cover of
  $\RP$, whose existence hinges on exactly the coherence
  $\ua(\mathsf{not}) \trns \ua(\mathsf{not}) \propeq \refl$,
  computes a $\mathbb{Z}/2$ winding and refutes $x \propeq \refl$.
  Together with $x \trns x \propeq \refl$ this machine-certifies
  that the internal realization presents $\mathbb{Z}/2$ faithfully:
  $x^2 = 1$ and $x \neq 1$ (\S\ref{sec:cover}).
\end{enumerate}

Along the way we report the measured computational frontier of the
kernel itself (\S\ref{sec:frontier}): transport along a single
$\hcomp$-composite of universe paths completes in about a minute,
while three nested layers blow up superlinearly --- the identified
obstruction, value-sharing across Kan steps, is a kernel-engineering
problem rather than a mathematical one, and the iterated per-letter
procedure sidesteps it entirely.

\subparagraph*{Relation to the companion papers.}  The kernel, the
constancy rule studied in \cite{Const}, and the cellular-context
platform of \cite{Sphere,Realization,Convexity} are reused here;
this paper is self-contained in its statements and cites the
companions for proofs of the imported facts
(\S\ref{sec:setting}).  The artifact conventions follow the series:
every machine-checked claim names a definition in the accompanying
repository (\S\ref{sec:artifact}).

\section{Setting and imported facts}
\label{sec:setting}

We work in the object theory of a defensively small cubical proof
kernel implementing a CCHM-style cubical type theory
\cite{CCHM}: dependent products and sums, universes
$\U_0 : \U_1$ (no type-in-type), a De Morgan interval $\I$ with
connections and reversal, $\Path$/$\mathsf{PathP}$ types,
$\transp$/$\hcomp$ with the standard computation rules, Glue types
with univalence ($\ua$) computing by unglueing, and the higher
inductive types used below ($\Sone$, pushouts).  Its algorithmic
equality $\aconv$ is a terminating normalization-by-evaluation
conversion check; \emph{decidability of $\aconv$ on well-typed
terms} is the kernel's basic metatheoretic budget, and everything
labeled ``machine-checked'' below is a type check (positive or, for
separations, deliberately failing) or a closed-term normalization
performed by this kernel.

\begin{defi}[cellular context]
  A \emph{cellular context} is a context of the object theory of
  the shape
  \[
    \Gamma \;=\; A : \U,\; a_1 : A,\; \dots,\;
    p_1 : \Path\,A\,a_i\,a_j,\; \dots,\;
    Q_1 : \mathsf{Sq}(\dots),\; \dots
  \]
  listing finitely many cell variables whose boundaries are formed
  from earlier cells --- the context of a finite cubical computad
  (polygraph).  A \emph{definable cell} over $\Gamma$ is a term of a
  cell type over $\Gamma$; a \emph{word} is a definable cell built
  from cell variables by path operations
  ($\refl$, inversion, $\cng$, composition $\trns$, and Kan
  fillers).
\end{defi}

We import three facts from the companions, with the anchor
statements re-proved in the artifact of this paper where they are
used.

\begin{enumerate}
\item[(K1)] \emph{Strict classification} \cite{Sphere}: over a
  generic cellular context, the strict inhabitants of a cell type
  are exactly the reparametrizations of cell variables along the
  free De Morgan algebra on the bound interval variables; the
  strict layer of a generic $n$-cell is the wedge-of-spheres
  monoid $\DM{-}$-set.
\item[(K2)] \emph{Realization} \cite{Realization}: cellular
  contexts realize their cube presheaves on the nose --- the
  attachment equations of the computad hold definitionally, and
  the kernel's $\aconv$ neither merges distinct presheaf cells nor
  misses an identification.
\item[(K3)] \emph{Decidability of the strict layer}
  \cite{Sphere,Realization}: equality of reparametrizations in the
  free De Morgan algebra is decidable (disjunctive normal forms of
  antichains), and $\aconv$ restricted to strict cells implements
  it.
\end{enumerate}

Two further companion results frame the interpretation: the no-go
theorems of \cite{Const,Convexity} (a nontrivial relator can never
be made definitional --- constant-tube $\hcomp$-regularity at path
types breaks canonicity-adjacent invariants), and the strict
contractibility of filler fibers \cite{Convexity}.  They will
reappear in \S\ref{sec:boundary} as the conceptual reason the
decidability boundary sits where it does.

\section{The strict word problem is decidable everywhere}
\label{sec:strict}

\begin{thm}
  \label{thm:strict}
  Over every cellular context $\Gamma$ and in every dimension, the
  strict word problem --- given two definable cells $u, v$ of a
  common cell type over $\Gamma$, decide $u \aconv v$ --- is
  decidable.  Moreover the decision agrees with equality in the cube
  presheaf presented by $\Gamma$: relator cells are treated as
  generating cells, never as equations.
\end{thm}

\begin{proof}
  Decidability is the kernel's conversion check on well-typed terms
  (\S\ref{sec:setting}).  Agreement with the presheaf semantics is
  (K2): the realization theorem states that $\aconv$ over $\Gamma$
  decides exactly the on-the-nose equality of the presented
  presheaf, with the strict layer classified by (K1) and its
  equational theory decided by (K3).  A relator cell --- a cell
  variable whose boundary happens to spell a word --- contributes
  generators to the presheaf, not equations: no computation rule of
  the kernel consults the boundary of a variable to rewrite a term.
  (This is the operational face of the no-go theorems: the kernel
  \emph{could not} consistently promote the relator to a rewrite
  even if asked; cf.\ \S\ref{sec:boundary}.)
\end{proof}

The theorem is deliberately unspectacular --- it records that the
strict layer stays decidable \emph{uniformly}, including over the
very contexts on which the propositional layer will turn
undecidable.  That uniformity is the left column of the table in
\S\ref{sec:intro}, and half of the boundary statement.

\section{Dimension one, free: the propositional problem is decidable}
\label{sec:free}

Let $\Gamma_G$ be a free graph context: a type $A$, finitely many
points, and finitely many path variables (edges), with no cells of
dimension $\geq 2$.  Words over $\Gamma_G$ are edge-words: composites
of edges and their inverses.  Classically, the free groupoid on a
graph has normal forms --- reduced words --- and two edge-words are
equal in the free groupoid iff their reductions coincide.  The
propositional layer of the kernel presents exactly the free
groupoid:

\begin{thm}
  \label{thm:free}
  Over a free graph context, the propositional word problem is
  decidable: $w_1 \propeq w_2$ is inhabited iff the reduced words of
  $w_1$ and $w_2$ coincide.  Both halves of the decision procedure
  are computational over the kernel:
  \begin{enumerate}
  \item (\emph{identification}) if the reductions coincide, an
    inhabitant of $w_1 \propeq w_2$ is compiled from the generic
    groupoid laws --- associativity, units, inverses ---
    each a machine-checked path over the generic context;
  \item (\emph{separation}) if the reductions differ, the words are
    separated by instantiating $\Gamma_G$ at a univalent cover and
    computing a transport invariant; since substitution preserves
    $\propeq$, distinct invariant values over the instance refute
    $w_1 \propeq w_2$ over the generic context.
  \end{enumerate}
\end{thm}

\begin{proof}[Proof scheme]
  Completeness of reduced words for the free groupoid is classical.
  (1) is a compilation: each reduction step (cancel $e \trns
  e^{-1}$, reassociate, absorb a unit) is an instance of a generic
  law available as a closed library path
  ($\mathsf{transAssoc}$, $\mathsf{transReflL/R}$,
  $\mathsf{transInvR/L}$ in the artifact), transported into word
  position by $\cng$ and chained by $\trns$; the compiled composite
  type-checks over $\Gamma_G$ itself.  (2) reduces to exhibiting,
  for each pair of distinct reduced words, an instance and an
  invariant with distinct computed values.  For a single edge
  instantiated at the circle, the winding invariant computes
  ($+1$, $-1$, $0$, $+2$, \dots{} are among the artifact's
  normalized values); for several edges the wedge of circles is the
  universal instance, and \S\ref{sec:sanov} supplies a complete,
  practically computable invariant.  Soundness of the refutation is
  transport invariance: an identification over $\Gamma_G$ would
  transport to an identification of distinct normal forms of the
  invariant's values.
\end{proof}

The interesting engineering content is in (2): completeness of the
separation half needs an invariant that (a) distinguishes all pairs
of distinct reduced words --- so abelian invariants are insufficient
--- and (b) actually evaluates on the kernel.  The next section
constructs one.

\section{A proof-free carrier: the Sanov invariant}
\label{sec:sanov}

The natural complete invariant for two generators is the
$\Ftwo$-cover of the figure eight $\Sone \vee \Sone$: fibers are
reduced words, monodromy prepends a generator with cancellation.
The series artifact contains this cover in full, machine-checked
--- and it is computationally impractical.  Its carrier is a
subset type (words with a reducedness proof), so every value carries
a proof, and the cover equivalence embeds a proof that the carrier
is a set; the kernel's Kan steps traverse those proofs.  A single
generator's winding normalizes in seconds; composites exhaust any
practical budget.  The failure mode is structural, not accidental:
\emph{proof-carrying carriers make transport walk proofs}.

The remedy is to change the invariant, not the kernel.  Sanov's
representation \cite{Sanov}
\[
  L \mapsto \begin{pmatrix}1 & 2\\ 0 & 1\end{pmatrix},
  \qquad
  R \mapsto \begin{pmatrix}1 & 0\\ 2 & 1\end{pmatrix}
\]
embeds $\Ftwo$ faithfully into $\SL$.  The carrier is now
$\mathbb{Z}^4$ --- a matrix, a pair of pairs of integers, with
\emph{no propositional payload}.

\begin{defi}[Sanov cover]
  \label{def:sanov}
  Right multiplication by $L$ (resp.\ $R$) is an equivalence of
  $\mathbb{Z}^4$: the inverse prepends the negated double, and both
  round trips reduce to two integer cancellation laws
  ($x + (-x + y) = y$ and its mirror), proved from associativity,
  inverses and units of $\mathbb{Z}$.  The \emph{Sanov cover}
  $\mathsf{helixSL} : \Sone \vee \Sone \to \U$ sends each circle to
  $\ua$ of its generator's equivalence; the \emph{matrix winding}
  $\mathsf{windSL}$ of a loop transports the identity matrix along
  the cover.
\end{defi}

\begin{thm}
  \label{thm:sanov}
  The Sanov invariant is complete for the free group on two
  generators and computes at machine speed on the kernel:
  \begin{enumerate}
  \item generators wind to their Sanov matrices in milliseconds
    (measured: $8\,$ms, against seconds for the reduced-word
    cover);
  \item the \emph{iterated} procedure --- one Glue transport per
    letter, feeding the concrete matrix forward --- runs at under a
    millisecond per letter; in particular
    $L\trns R \mapsto \left(\begin{smallmatrix}5&2\\2&1
    \end{smallmatrix}\right)$ and
    $R\trns L \mapsto \left(\begin{smallmatrix}1&2\\2&5
    \end{smallmatrix}\right)$ are distinct --- separating a pair
    that every abelian invariant conflates --- and the commutator
    $L\trns R\trns L^{-1}\trns R^{-1}$ evaluates to
    $\left(\begin{smallmatrix}21&-8\\8&-3\end{smallmatrix}\right)
    \neq I$;
  \item two words are equal in $\Ftwo$ iff their computed matrices
    coincide (faithfulness of the representation), so the invariant
    decides the separation half of Theorem~\ref{thm:free}
    completely, in time linear in word length.
  \end{enumerate}
  The agreement of the iterated procedure with the direct winding of
  the composite loop is propositional and generic: $\cng$
  distributes over composition
  ($\mathsf{congTrans}$), and transport along a composite of
  universe paths is the composite of transports
  ($\mathsf{transpTrans}$), both closed library paths.
\end{thm}

\begin{proof}
  (1) and (2) are kernel computations, machine-checked against the
  expected matrices in the artifact's native harness; (3) is
  Sanov's theorem \cite{Sanov} together with soundness of the
  computed values (canonicity for the closed integer type: a
  normalized matrix entry is a numeral, and distinct numerals are
  separated by the decidable equality of $\mathbb{Z}$, itself a
  library fact).  The bridging identifications are the named
  library paths.
\end{proof}

\begin{rem}
  The lesson we draw --- and use again in \S\ref{sec:cover} --- is a
  design principle for computational univalence:
  \emph{represent invariants on proof-free carriers, pushing all
  proof obligations into the equivalences}.  The equivalence proofs
  are walked once per Glue formation, not once per transported
  value; normal forms of values stay small; and a classical
  faithful representation (here into $\SL$) converts completeness of
  the invariant into a theorem about the target group rather than
  about the type-theoretic carrier.
\end{rem}

\section{Relators compile: the presentation context}
\label{sec:relators}

From dimension~$2$ on, cellular contexts contain \emph{relator
cells}: squares whose boundary spells a word.  Up to the standard
globular--cubical translation, the presentation context of
$\langle x \mid x^2 \rangle$ is
\[
  \Gamma_{\mathbb{Z}/2}
  \;=\;
  A : \U,\;
  a : A,\;
  x : \Path\,A\,a\,a,\;
  R : \Path\,(\Path\,A\,a\,a)\,(x \trns x)\,\refl .
\]
Over $\Gamma_{\mathbb{Z}/2}$, the group-theoretic consequences of
the relation compile to explicit paths, exactly as the groupoid
laws did in Theorem~\ref{thm:free}(1) --- with $R$ now among the
usable moves:

\begin{prop}
  \label{prop:compile}
  Over $\Gamma_{\mathbb{Z}/2}$ there is a machine-checked path
  \begin{gather*}
    \mathsf{wordX4}
    \;:\;
    (x \trns x) \trns (x \trns x) \;\propeq\; \refl,\\
    \mathsf{wordX4}
    \;=\;
    \cng_{(-)\trns(x\trns x)}(R)
    \trns \mathsf{transReflL}
    \trns R ,
  \end{gather*}
  i.e.\ $x^4 = 1$ is derived by two relator uses and one unit law.
  The compilation scheme is uniform: a derivation of $w = 1$ in the
  presented group maps to a path $w \propeq \refl$ by sending each
  relator application to a $\cng$-embedded use of the relator
  variable and each free-reduction step to a groupoid law, chained
  by $\trns$.
\end{prop}

This is one half of the undecidability transfer: \emph{provable
group equalities become inhabitants}.  It is worth noting what the
compilation does \emph{not} need: no instance, no Kan machinery
beyond composition, no knowledge of the ambient type $A$ --- the
derivation lives over the generic context and transfers to every
model by substitution.  The other half --- \emph{inhabitants imply
group equalities} --- needs an interpretation, hence a realization
of the presentation complex; that is the next section.

\section{Presentations realize: $\RP$ inside the kernel}
\label{sec:present}

The presentation complex of $\langle x \mid x^2 \rangle$ --- one
vertex, one edge, one $2$-cell attached along $x^2$ --- is the
mapping cone of the degree-$2$ self-map of the circle: the real
projective plane.  Both the attaching map and the cone exist as
first-class objects of the kernel:

\begin{defi}
  \label{def:rp2}
  $\deg_2 : \Sone \to \Sone$ is the circle eliminator with
  $\mathsf{base} \mapsto \mathsf{base}$ and $\mathsf{loop} \mapsto
  \mathsf{loop} \trns \mathsf{loop}$;
  $\RP := \mathsf{cofib}(\deg_2)$ is its mapping cone, the pushout
  $\mathbf{1} \leftarrow \Sone \xrightarrow{\deg_2} \Sone$.  Write
  $b := \mathsf{inr}(\mathsf{base})$ and
  $x := \langle i \rangle\, \mathsf{inr}(\mathsf{loop}\,i)$ for the
  realized vertex and generator.
\end{defi}

\begin{thm}
  \label{thm:realize}
  The internal $\RP$ realizes the presentation:
  \begin{enumerate}
  \item (\emph{attachment, strictly}) $\cng_{\deg_2}(\mathsf{loop})
    \aconv \mathsf{loop} \trns \mathsf{loop}$ --- the attaching map
    computes definitionally on the generator ---
    and the cone path from the cone point to $b$ is a definable
    cell;
  \item (\emph{the relator cell, realized}) there is a
    machine-checked identification
    $\mathsf{rp2Rel} : x \trns x \propeq \refl_b$, obtained by
    folding the pushout cylinder with a single $\hcomp$ square
    capped at the cone point, then splitting the image of
    $\mathsf{loop}\trns\mathsf{loop}$ by $\mathsf{congTrans}$;
  \item (\emph{transfer by instantiation}) the substitution of
    $(\RP, b, x, \mathsf{rp2Rel})$ for
    $(A, a, x, R)$ in Proposition~\ref{prop:compile} yields a
    machine-checked identification $x^4 \propeq \refl$ in $\RP$:
    generic relator derivations land on the internal realization
    with no further work;
  \item (\emph{properly propositional}) the identification in (2)
    is not strict: the constant path-lambda fails to type-check
    against $x \trns x \propeq \refl$ (a deliberate negative probe),
    as the no-go theorems require.
  \end{enumerate}
\end{thm}

\begin{proof}
  All four items are single artifact probes.  (1) the constant
  path-lambda inhabits
  $\langle i\rangle \deg_2(\mathsf{loop}\,i) \propeq
  \mathsf{loop}\trns\mathsf{loop}$, so the attachment is
  definitional; (2) the square has the pushout cylinder as its
  $k{=}0$ wall, the cone point as cap, and the basepoint cone path
  on the remaining three walls; (3) is literally an application of
  the closed term $\mathsf{wordX4}$; (4) is a failing type check,
  recorded as a guard.
\end{proof}

\section{Faithfulness: the double cover of $\RP$}
\label{sec:cover}

Theorem~\ref{thm:realize} would be vacuous if the realization
collapsed --- if, say, $x \propeq \refl$ held in $\RP$, every word
would be identified and the relator would certify nothing.  The
classical statement $\pi_1(\RP) = \mathbb{Z}/2$ says precisely that
the presentation is realized \emph{faithfully}: $x^2 = 1$ (the
relator) and $x \neq 1$ (no collapse).  We certify the computational
core of this on the kernel with a univalent double cover, built by
the proof-free-carrier principle of \S\ref{sec:sanov} with carrier
$\mathsf{Bool}$.

\begin{lem}
  \label{lem:uanotnot}
  $\ua(\mathsf{not}) \trns \ua(\mathsf{not}) \propeq
  \refl_{\mathsf{Bool}}$, machine-checked (by injectivity of $\ua$,
  equality of equivalences from equality of functions, and the
  pointwise involution $\mathsf{not}\circ\mathsf{not} = \mathsf{id}$).
\end{lem}

\begin{thm}
  \label{thm:cover}
  There is a cover $\mathsf{covRP2} : \RP \to \U$ with fiber
  $\mathsf{Bool}$, sending the $1$-skeleton circle to
  $\ua(\mathsf{not})$; its extension over the relator $2$-cell
  requires exactly Lemma~\ref{lem:uanotnot}, and is constructed by
  transporting the trivial square along it (path induction in
  $\U_1$).  The resulting winding
  $\mathsf{windRP2} := \transp_{\mathsf{covRP2}}(\mathsf{true})$
  computes:
  \[
    \mathsf{windRP2}(\refl) \longmapsto \mathsf{true},
    \qquad
    \mathsf{windRP2}(x) \longmapsto \mathsf{false},
  \]
  whence a machine-checked refutation
  \[
    \mathsf{rp2LoopNontriv} \;:\;
    (x \propeq \refl) \to \mathbf{0}.
  \]
  Together with $\mathsf{rp2Rel}$
  (Theorem~\ref{thm:realize}(2)), the presentation
  $\langle x \mid x^2\rangle$ is realized faithfully at the level of
  loop identifications: $x^2 = 1$ and $x \neq 1$.  As a corollary,
  $\RP$ is not a set.
\end{thm}

\begin{proof}
  The cover is a pushout eliminator into $\U$; on the cone it is
  constantly $\mathsf{Bool}$, on the $1$-skeleton it is the circle
  eliminator with loop payload $\ua(\mathsf{not})$, and on the
  pushout cylinder it must provide, over each point of the attaching
  circle, a path $\mathsf{Bool} \propeq
  \mathsf{cov}(\deg_2(-))$ --- a square whose monodromy side is the
  image of $\mathsf{loop} \trns \mathsf{loop}$, i.e.\ (after
  $\mathsf{congTrans}$) the composite
  $\ua(\mathsf{not}) \trns \ua(\mathsf{not})$.
  Lemma~\ref{lem:uanotnot} is thus literally the extension
  condition; path induction along it reduces the required square to
  the trivial one.  The two windings are closed-term
  normalizations; the refutation composes
  $\cng_{\mathsf{windRP2}}$ with the decidable separation of
  $\mathsf{false}$ from $\mathsf{true}$.
\end{proof}

\begin{rem}
  The cover makes a second architectural point: the extension
  condition over a relator cell is always a coherence in the
  universe --- here $\ua(\mathsf{not})^{\trns 2} \propeq \refl$,
  in general the image of the relator word under the intended
  monodromy.  Constructing covers of internal presentation
  complexes is thus exactly as hard as proving the corresponding
  relation among $\ua$-paths, a task the proof-free-carrier
  principle keeps tractable.
\end{rem}

\section{From dimension two: undecidability}
\label{sec:undec}

\begin{thm}
  \label{thm:undec}
  The propositional word problem over cellular contexts with cells
  of dimension $\geq 2$ is undecidable: there is no algorithm that,
  given a cellular context $\Gamma$ and two edge-words $w_1, w_2$
  over it, decides whether $w_1 \propeq w_2$ is inhabited.
\end{thm}

\begin{proof}[Proof, by reduction]
  Let $\langle X \mid \mathcal{R} \rangle$ be a finite group
  presentation and $\Gamma_{\langle X \mid \mathcal{R}\rangle}$ its
  presentation context: one point, one loop per generator, one
  relator square per relation (as in \S\ref{sec:relators}).  We
  claim $w \propeq \refl$ is inhabited over
  $\Gamma_{\langle X \mid \mathcal{R}\rangle}$ iff $w = 1$ in the
  presented group $G$; undecidability then follows from the
  existence of finitely presented groups with undecidable word
  problem \cite{Novikov,Boone}.

  ($\Leftarrow$) is the compilation of
  Proposition~\ref{prop:compile}, verbatim for arbitrary
  presentations.  ($\Rightarrow$): interpret
  $\Gamma_{\langle X \mid \mathcal{R}\rangle}$ into the realization
  of its presentation complex $K_{\langle X \mid \mathcal{R}\rangle}$
  --- the interpretation exists in any model of the theory with
  pushouts, sending cell variables to the corresponding cells, as in
  Theorem~\ref{thm:realize} for the one-relator case --- and apply
  $\pi_1(K_{\langle X \mid \mathcal{R}\rangle}) = G$ (van Kampen):
  an inhabitant of $w \propeq \refl$ over the generic context
  substitutes to an identification in the realization, hence to
  $w = 1$ in $G$.  Soundness of the type theory for such models is
  the standard cubical-sets semantics \cite{CCHM}.
\end{proof}

\begin{rem}
  The kernel-checked material of
  \S\S\ref{sec:relators}--\ref{sec:cover} exhibits both directions
  of the reduction in the minimal nontrivial instance
  $\mathbb{Z}/2$: the compilation ($\Leftarrow$) as
  $\mathsf{wordX4}$, and the faithful realization underlying
  ($\Rightarrow$) as
  $\mathsf{rp2Rel}$ + $\mathsf{rp2LoopNontriv}$ ---
  including the genuinely propositional character of the relator
  (Theorem~\ref{thm:realize}(4)).  We do not claim an internal
  proof of $\pi_1 = G$ for arbitrary presentations; the reduction
  above uses it classically, as is standard for undecidability
  transfers.
\end{rem}

\begin{rem}
  A cognate undecidability result was obtained independently by
  Dor\'e, Cavallo and M\"ortberg \cite[Thm.~3.23]{DCM}: filling a
  given \emph{Kan boundary} by a Kan composite, over the same
  interval algebras, is undecidable, again by reduction to group
  word problems.  The two statements sit on either side of this
  paper's boundary table: theirs concerns the existence of a single
  Kan filler for a fixed boundary (a question about the weak layer's
  operations), ours the propositional equality of words over relator
  contexts; on the strict side, both their contortion solving and
  our Theorem~\ref{thm:strict} remain decidable.
\end{rem}

\section{The boundary}
\label{sec:boundary}

Assembling Theorems~\ref{thm:strict}, \ref{thm:free}
and~\ref{thm:undec}:

\begin{center}
\begin{tabular}{l|l|l}
  & dimension $1$ (free) & dimension $\geq 2$ (relators)\\
  \hline
  strict $\aconv$
    & decidable (K1--K3)
    & decidable (Theorem~\ref{thm:strict})\\
  weak $\propeq$
    & decidable (Theorem~\ref{thm:free})
    & undecidable (Theorem~\ref{thm:undec})
\end{tabular}
\end{center}

The two layers separate exactly at the relator cells, and the no-go
theorems of the series say this is not an accident of
implementation.  A kernel that tried to decide the propositional
layer by promoting relators to definitional equations would need
$R : x\trns x \propeq \refl$ to hold strictly; the no-go theorems
\cite{Const,Convexity} show that even the humblest such promotion
--- constant-tube $\hcomp$-regularity at path types --- already
destroys the invariants behind canonicity.  Undecidability gives the
quantitative version: \emph{no} extension of the strict equality
can absorb the relators of an arbitrary presentation while remaining
decidable, since the resulting theory would decide group word
problems.  Conversely, the price is worth paying exactly where it is
paid: the weak layer's identifications are what make the
presentation complexes, covers and invariants of
\S\S\ref{sec:present}--\ref{sec:cover} exist at all.  The
strict--weak boundary is not a defect line but a load-bearing wall:
on one side, a decidable combinatorics of reparametrizations
(K1--K3); on the other, the full undecidable expressiveness of
homotopy theory --- with the compilation of
Proposition~\ref{prop:compile} and the covers of
\S\ref{sec:cover} as the traffic across it.

\section{The measured frontier}
\label{sec:frontier}

Two performance walls were identified and measured in the course of
this work; both concern the kernel, not the mathematics, and both
are reported here as data for kernel engineering.

\emph{Proof-carrying carriers} (\S\ref{sec:sanov}): transport walks
the proofs embedded in carrier values and cover equivalences.  With
the reduced-word carrier a single generator costs seconds and
composites are infeasible; the proof-free Sanov carrier removes this
wall entirely (milliseconds per letter, measured).

\emph{Nested Kan composition in the universe}: transporting
directly along a composite loop evaluates $\hcomp$ in $\U$ over
Glue lines.  Measured on the Sanov cover: one nested layer (the
cancelling composite, or the wedge-conjugated second generator)
completes in about a minute; three nested layers exceed 50 CPU
minutes without completing.  The growth is superlinear in nesting
depth; the identified cause is re-evaluation of inner Glue layers at
each Kan step, and the identified remedy --- value-sharing
(memoization) across Kan steps --- is a kernel-engineering project
orthogonal to this paper.  The iterated per-letter procedure of
Theorem~\ref{thm:sanov} avoids the wall by construction, and its
agreement with the direct composite is the generic propositional
bridge ($\mathsf{congTrans}$, $\mathsf{transpTrans}$); nothing in
Theorems~\ref{thm:free}--\ref{thm:undec} depends on the wall.

\section{Outlook}
\label{sec:outlook}

Three continuations suggest themselves.  First, \emph{internal
$\pi_1$ for general presentations}: the $\mathbb{Z}/2$ case of
\S\S\ref{sec:present}--\ref{sec:cover} used no property of the
presentation beyond its finiteness; an internal encode--decode proof
of $\pi_1(K_{\langle X \mid \mathcal{R} \rangle}) = G$ (in the style
of \cite{LicataShulman}, with the cover extension conditions
discharged relator-by-relator as in Lemma~\ref{lem:uanotnot}) would
make the ($\Rightarrow$) direction of Theorem~\ref{thm:undec}
internal as well.  Second, \emph{up-to-homotopy normal forms}: over
free contexts, reduced words are normal forms for the weak layer;
whether a useful normal-form calculus survives some relator classes
(one-relator groups, small cancellation) is the type-theoretic
shadow of classical combinatorial group theory, with the compiled
identifications of Proposition~\ref{prop:compile} as the rewriting
moves.  Third, \emph{the kernel frontier}: value-sharing across Kan
steps (\S\ref{sec:frontier}) would make direct composite transport
practical and is the single most valuable kernel improvement
identified by this series.

\subparagraph*{Acknowledgements.}  The platform and companion
results were developed in \cite{Const,Sphere,Realization,Convexity}.

\section{Artifact}
\label{sec:artifact}

\sloppy
All machine-checked claims are definitions with positive or negative
guards in the accompanying repository (kernel, library and probe
suites; the native harness reports the timings quoted above).
Anchors: the presentation-context compilation
(\S\ref{sec:relators}) is \texttt{wordX4}/\texttt{wordRel}; the
Sanov invariant (\S\ref{sec:sanov}) is
\texttt{addCancel}/\texttt{addCancelN}, \texttt{sanovL/R},
\texttt{helixSL}, \texttt{windSL}, with the timing harness
\texttt{Test/Sanov.lean}; the internal presentation complex
(\S\ref{sec:present}) is \texttt{deg2}, \texttt{deg2Loop},
\texttt{rp2}, \texttt{rp2Attach}, \texttt{rp2RelSq},
\texttt{rp2Rel}, \texttt{wordX4RP2}, \texttt{wordRelRP2}; the double
cover (\S\ref{sec:cover}) is \texttt{uaNotNot}, \texttt{covS1},
\texttt{covRP2}, \texttt{windRP2}, \texttt{rp2LoopNontriv},
\texttt{rp2NotSet}; the groupoid laws and winding values of
Theorem~\ref{thm:free} and the reduced-word $\Ftwo$ cover are in the
library of the earlier releases.  Repository:
\url{https://github.com/r-shibusawa/cubical-strict-j}, archived at
DOI \href{https://doi.org/10.5281/zenodo.21405961}%
{10.5281/zenodo.21405961} (concept DOI; the release accompanying
this paper is tagged in the repository).

\begin{thebibliography}{10}

\bibitem{Boone}
W.~W.~Boone.
\newblock The word problem.
\newblock \emph{Ann.\ of Math.}\ 70 (1959), 207--265.

\bibitem{CCHM}
C.~Cohen, T.~Coquand, S.~Huber, A.~M\"ortberg.
\newblock Cubical type theory: a constructive interpretation of the
  univalence axiom.
\newblock \emph{TYPES 2015}, LIPIcs 69, 2018.

\bibitem{DCM}
M.~Dor\'e, E.~Cavallo, A.~M\"ortberg.
\newblock Automating boundary filling in cubical type theories.
\newblock \emph{Log.\ Methods Comput.\ Sci.}\ 22(2):28, 2026.

\bibitem{LicataShulman}
D.~R.~Licata, M.~Shulman.
\newblock Calculating the fundamental group of the circle in
  homotopy type theory.
\newblock \emph{LICS 2013}.

\bibitem{Novikov}
P.~S.~Novikov.
\newblock On the algorithmic unsolvability of the word problem in
  group theory.
\newblock \emph{Trudy Mat.\ Inst.\ Steklov}\ 44 (1955), 1--143.

\bibitem{Sanov}
I.~N.~Sanov.
\newblock A property of a representation of a free group.
\newblock \emph{Dokl.\ Akad.\ Nauk SSSR}\ 57 (1947), 657--659.

\bibitem{Const}
R.~Shibusawa.
\newblock Evaluation-time constancy inference for transport in
  cubical type theory.
\newblock Preprint, 2026.  Artifact:
  \url{https://github.com/r-shibusawa/cubical-strict-j}, archived
  at DOI \href{https://doi.org/10.5281/zenodo.21637898}%
  {10.5281/zenodo.21637898}.

\bibitem{Sphere}
R.~Shibusawa.
\newblock The strict layer of generic cells in cubical type theory
  is a wedge of De Morgan spheres.
\newblock Preprint, 2026.  Included in the artifact repository
  above (from release \texttt{v1.3.0}, DOI
  \href{https://doi.org/10.5281/zenodo.21638976}%
  {10.5281/zenodo.21638976}).

\bibitem{Realization}
R.~Shibusawa.
\newblock Generic cellular contexts in cubical type theory realize
  their cube presheaves on the nose.
\newblock Preprint, 2026.  Included in the artifact repository
  above (from release \texttt{v1.3.0}, DOI as for \cite{Sphere}).

\bibitem{Convexity}
R.~Shibusawa.
\newblock De Morgan convexity and the strict--weak comparison in
  cubical type theory.
\newblock Preprint, 2026.  Included in the artifact repository
  above (from release \texttt{v1.4.0}, DOI
  \href{https://doi.org/10.5281/zenodo.21639454}%
  {10.5281/zenodo.21639454}).

\end{thebibliography}

\end{document}
